\documentclass[11pt]{article}

\usepackage[margin=1in]{geometry}
\usepackage{listings}
\usepackage{booktabs}
\usepackage[colorlinks=true,linkcolor=blue,citecolor=blue,urlcolor=blue]{hyperref}
\usepackage{xcolor}
\usepackage{graphicx}
\usepackage{microtype}
\usepackage{url}
\usepackage{etoolbox}

%%% ---- Language definitions for listings ----

\lstdefinelanguage{Modula3}{
  morekeywords={MODULE,INTERFACE,IMPORT,FROM,EXPORT,TYPE,CONST,VAR,
    PROCEDURE,BEGIN,END,IF,THEN,ELSIF,ELSE,WHILE,DO,FOR,TO,BY,
    REPEAT,UNTIL,LOOP,EXIT,WITH,CASE,OF,RETURN,RAISE,RAISES,
    TRY,EXCEPT,FINALLY,LOCK,NEW,NARROW,ISTYPE,TYPECASE,TYPECODE,
    BRANDED,REF,REFANY,OBJECT,METHODS,OVERRIDES,REVEAL,RECORD,
    ARRAY,SET,NIL,TRUE,FALSE,NOT,AND,OR,IN,MOD,DIV,ABS,
    FIRST,LAST,NUMBER,ORD,VAL,FLOAT,CEILING,FLOOR,ROUND,TRUNC,
    ADR,LOOPHOLE,EVAL,EXTERNAL,UNUSED,CARDINAL,INTEGER,LONGINT,
    REAL,LONGREAL,EXTENDED,BOOLEAN,CHAR,TEXT},
  sensitive=true,
  morecomment=[s]{(*}{*)},
  morestring=[d]{"},
}

\lstdefinelanguage{Scheme}{
  morekeywords={define,lambda,let,let*,letrec,if,cond,else,and,or,
    not,begin,set!,quote,import,car,cdr,cons,list,map,filter,
    for-each,apply,eval,display,newline,eq?,eqv?,equal?,null?,
    pair?,number?,string?,symbol?,boolean?,char?,procedure?},
  sensitive=true,
  morecomment=[l]{;},
  morestring=[d]{"},
  alsoother={!,?},
}

\lstset{
  basicstyle=\small\ttfamily,
  breaklines=true,
  frame=single,
  xleftmargin=2em,
  framexleftmargin=1.5em,
  numbers=none,
  showstringspaces=false,
  columns=flexible,
  keepspaces=true,
  aboveskip=1em,
  belowskip=1em,
}

% Prevent code listings and verbatim from breaking across pages
\BeforeBeginEnvironment{lstlisting}{\begin{minipage}{\linewidth}}
\AfterEndEnvironment{lstlisting}{\end{minipage}}

%%% ---- Document ----

\title{MScheme: Modula-3/Scheme Interoperability\\and Comparable Systems}
\author{Mika Nystr\"om \and Claude (Anthropic)}
\date{March 5, 2026}

\begin{document}

\maketitle

\begin{abstract}
MScheme is a Scheme interpreter written in Modula-3.  Because
Modula-3's universal traced reference type (\texttt{REFANY}) serves
as the Scheme value type, every host-language object is already a
valid Scheme value---no wrapping or boxing is needed.  A build-time
stub generator reads standard Modula-3 interfaces and produces
bidirectional bindings automatically.  We describe the architecture,
compare it to ten other embedding systems, and place it in historical
context.
\end{abstract}

\tableofcontents

%% ==================================================================
\section{Introduction}
\label{sec:intro}
%% ==================================================================

MScheme is a Scheme interpreter written in Modula-3, shipped with a
stub generator (\texttt{sstubgen}) that processes Modula-3 interface
files and produces bidirectional bindings.  It originated as a
line-for-line translation of Peter Norvig's JScheme
(Java)~\cite{norvig1998}.  The key insight that makes JScheme's
Java/Scheme interoperability seamless is that if the host language has
a universal reference type, then every host-language value already
\emph{is} a valid Scheme value---no wrapping or boxing is needed.  In
Java that type is \texttt{java.lang.Object}; in Modula-3 it is
\texttt{REFANY}.  MScheme carries this principle into the Modula-3
world.

This report describes MScheme's interoperability architecture in
detail, compares it to other embedding systems that solve similar
problems in different ways, and places the design in historical
context---from its direct predecessor Tinylisp (DEC SRC,
1989)~\cite{ellis1989} through the Lisp Machines to the broader
question of how Algol-family languages can support this architecture
on conventional hardware.

\subsection*{A taste of MScheme}

Before diving into the architecture, here is real code from
\href{https://github.com/IntelLabs/async-toolkit/tree/main/m3utils/m3utils/csp}{\texttt{cspc}},
the CSP compiler built on MScheme.  This is Scheme code running
inside a Modula-3 program:

\begin{lstlisting}[language=Scheme]
;; Read an environment variable -- Env is an M3 module
(define *m3utils* (Env.Get "M3UTILS"))

;; Create M3 BigInt objects -- BigInt.T <: REFANY, so they are
;; ordinary Scheme values
(define a (BigInt.New 12))
(define b (BigInt.New 42))
(BigInt.Equal a b)          ;=> #f
(BigInt.Format a 10)        ;=> "12"

;; M3 procedures used as first-class Scheme values:
;; BigInt.Equal and BigInt.Compare are passed directly to lambda
(define *boolean-binops*
  (list
    `(== integer ,BigInt.Equal)
    `(<  integer ,(lambda (a b) (< (BigInt.Compare a b) 0)))
    `(>  integer ,(lambda (a b) (> (BigInt.Compare a b) 0)))))

;; Format an integer literal by calling M3's Fmt.Int
(define (format-int-literal x)
  (Fmt.Int (BigInt.ToInteger x) 10))

;; Create an M3 object, call its methods from Scheme
(define the-driver
  (obj-method-wrap
    (new-modula-object 'CspCompilerDriver.T)
    'CspCompilerDriver.T))
(the-driver 'init "hello.procs")
(the-driver 'getProcTypes)  ;=> a TextSeq.T (M3 sequence object)
\end{lstlisting}

Every value here---the \texttt{BigInt.T} integers, the \texttt{TEXT}
strings, the \texttt{CspCompilerDriver.T} object---is a native
Modula-3 heap object.  None of them are wrapped, copied, or
marshaled.  They are Scheme values because every M3 reference type
already satisfies \texttt{SchemeObject.T = REFANY}.  The procedure
names (\texttt{BigInt.Equal}, \texttt{Fmt.Int}, etc.)\ were generated
automatically by \texttt{sstubgen} from the corresponding Modula-3
interface files at build time.  The generated stub files live in the
platform-specific build directory (e.g.,
\texttt{csp/ARM64\_DARWIN/FmtSchemeStubs.m3}).


%% ==================================================================
\section{MScheme Architecture}
\label{sec:architecture}
%% ==================================================================

\subsection{The Universal Value Type}
\label{sec:universal-value}

The foundation of the entire system is one line of Modula-3:

\begin{lstlisting}[language=Modula3]
(* SchemeObject.i3 *)
TYPE T = REFANY;
\end{lstlisting}

A Scheme value is \texttt{REFANY}---Modula-3's traced, polymorphic
reference type.  Any heap-allocated M3 value (objects, REFs, TEXTs,
arrays) already \emph{is} a valid Scheme object.  No wrapper structs,
tag bits on the Scheme side, or special allocation paths.  The M3
runtime's garbage collector manages all values uniformly.

Scheme-specific types are defined as ordinary M3 reference types:

\begin{table}[ht]
\centering
\begin{tabular}{@{}ll@{}}
\toprule
Scheme type  & Modula-3 representation \\
\midrule
\multicolumn{2}{@{}l}{\emph{Numeric tower (exact $\to$ inexact $\to$ compound):}} \\
\quad Exact integer  & \texttt{SchemeInt.T = REF INTEGER} (fixnum, cached $[-256..256]$) \\
\quad               & \texttt{Mpz.T} (bignum via GMP) \\
\quad Exact rational & \texttt{SchemeRational.T} (num/den as \texttt{Mpz.T}) \\
\quad Inexact real   & \texttt{SchemeLongReal.T = REF LONGREAL} \\
\quad Arb-prec float & \texttt{SchemeMpfr.T} (MPFR wrapper, configurable precision) \\
\quad Complex        & \texttt{SchemeComplex.T} (re/im as \texttt{SchemeObject.T}) \\
\quad Dual           & \texttt{SchemeDual.T} (re/$\varepsilon$ for automatic differentiation) \\
Pair         & \texttt{SchemePair.T = REF RECORD first, rest : REFANY END} \\
Symbol       & \texttt{SchemeSymbol.T = Atom.T} \\
Boolean      & \texttt{SchemeBoolean.T = BRANDED REF BOOLEAN} \\
String       & \texttt{SchemeString.T = REF ARRAY OF CHAR} \\
Procedure    & \texttt{SchemeProcedure.T} (object with \texttt{apply} method) \\
Any M3 ref   & Itself---already \texttt{REFANY} \\
\bottomrule
\end{tabular}
\caption{Scheme-specific types as Modula-3 reference types.}
\label{tab:scheme-types}
\end{table}

Runtime type discrimination uses Modula-3's \texttt{TYPECASE} (a
type-safe downcast), \texttt{ISTYPE}, and \texttt{TYPECODE}---the
same mechanisms the M3 runtime provides for all M3 code.  No parallel
type system is needed.

Modula-3 distinguishes \emph{reference types} (heap-allocated, traced
by the GC) from \emph{value types} (unboxed scalars, records,
arrays).  Reference types pass through the Scheme boundary unchanged,
but value types must be boxed into reference types.  The
\texttt{sstubgen}-generated stubs handle this automatically:

\begin{table}[ht]
\centering
\small
\begin{tabular}{@{}lll@{}}
\toprule
M3 value type & Scheme representation & Conversion \\
\midrule
\texttt{INTEGER}     & \texttt{SchemeInt.T}      & \texttt{SchemeInt.FromI} / \texttt{ToModula\_INTEGER} \\
\texttt{CARDINAL}    & \texttt{SchemeInt.T}      & Range-checked \\
\texttt{LONGREAL}    & \texttt{SchemeLongReal.T} & \texttt{FromLR} / \texttt{FromO} \\
\texttt{REAL}        & \texttt{SchemeLongReal.T} & Promoted via \texttt{FLOAT} \\
\texttt{EXTENDED}    & \texttt{SchemeLongReal.T} & Promoted via \texttt{FLOAT} \\
\texttt{Mpz.T}      & \texttt{SchemeInt.T} or \texttt{Mpz.T} & \texttt{SchemeInt.MpzToScheme} \\
\texttt{Mpfr.T}     & \texttt{SchemeMpfr.T}     & Preserves precision \\
\texttt{BOOLEAN}     & \texttt{SchemeBoolean.T}  & \texttt{Truth} / \texttt{TruthO} \\
\texttt{CHAR}        & \texttt{SchemeChar.T} (\texttt{REF CHAR}) & Cached; all 256 values preallocated \\
\texttt{TEXT}        & \texttt{SchemeString.T}   & \texttt{FromText} / \texttt{ToText} \\
Enumeration          & \texttt{SchemeInt.T}      & As ordinal index (0, 1, 2, \ldots) \\
Subrange             & \texttt{SchemeInt.T}      & As base type, with bounds check \\
\bottomrule
\end{tabular}
\caption{Value type conversions generated by \texttt{sstubgen}.}
\label{tab:value-types}
\end{table}

The numeric tower preserves exactness: exact integers (\texttt{SchemeInt.T},
\texttt{Mpz.T}) and exact rationals (\texttt{SchemeRational.T}) remain exact;
inexact reals use \texttt{SchemeLongReal.T} or \texttt{SchemeMpfr.T};
compound types \texttt{SchemeComplex.T} and \texttt{SchemeDual.T} compose
over the tower.  Common fixnum values ($[-256..256]$) and common
\texttt{LONGREAL} values (0.0, 1.0, $-$1.0, 2.0) are cached to reduce
allocation.  The \texttt{SchemeChar} module preallocates all
256 character values at initialization.

Records and arrays require more elaborate treatment because they are
\emph{aggregate} value types.  In Modula-3, a \texttt{RECORD} is a
stack-allocated or field-embedded aggregate (like a C struct) and an
\texttt{ARRAY} is a contiguous block of elements with fixed bounds.
Neither is a reference---neither can be \texttt{REFANY}.  Rather than
boxing them into opaque \texttt{REF} wrappers (which would make the
contents inaccessible to Scheme), \texttt{sstubgen} \emph{decomposes}
them into Scheme data structures:

\textbf{Records become association lists.}  Each field becomes a
\texttt{(name . value)} pair, with the field value converted
recursively.  A record \texttt{RECORD x: INTEGER; y: REAL END} with
values \texttt{x=3, y=1.5} becomes \texttt{((x . 3.0) (y . 1.5))}.
The generated \texttt{ToScheme} converter iterates over fields; the
\texttt{ToModula} converter walks the alist, matching field names by
symbol identity and assigning back into a stack-allocated M3 record:

\begin{lstlisting}[language=Modula3]
(* Generated by sstubgen for RECORD x: INTEGER; y: REAL END *)
PROCEDURE ToModula_MyRecord(x: SchemeObject.T): MyRecord =
  VAR res: MyRecord; p := SchemePair.Pair(x);
  BEGIN
    WHILE p # NIL DO
      IF    sym_x = First(p.first) THEN res.x := Int(Rest(p.first))
      ELSIF sym_y = First(p.first) THEN res.y := Float(Rest(p.first))
      END;
      p := SchemePair.Pair(p.rest)
    END;
    RETURN res
  END ToModula_MyRecord;
\end{lstlisting}

\textbf{Fixed arrays become lists of (index, element) pairs.}  An
\texttt{ARRAY [0..2] OF INTEGER} with values \texttt{\{10, 20, 30\}}
becomes \texttt{((0 . 10.0) (1 . 20.0) (2 . 30.0))}.  The
\texttt{ToModula} converter also accepts a simplified form without
explicit indices---a plain list \texttt{(10.0 20.0 30.0)}---assigning
elements to successive indices from \texttt{FIRST} to \texttt{LAST}.
The generated converter handles both forms:

\begin{lstlisting}[language=Modula3]
(* Generated by sstubgen for ARRAY [0..2] OF INTEGER *)
PROCEDURE ToModula_A(x: SchemeObject.T): ARRAY [0..2] OF INTEGER =
  VAR res: ARRAY [0..2] OF INTEGER; p := SchemePair.Pair(x);
  BEGIN
    IF p # NIL AND NOT ISTYPE(p.first, SchemePair.T) THEN
      (* Simple form: (10 20 30) *)
      FOR i := FIRST(res) TO LAST(res) DO
        res[i] := Int(p.first); p := SchemePair.Pair(p.rest)
      END
    ELSE
      (* Indexed form: ((0 . 10) (1 . 20) (2 . 30)) *)
      WHILE p # NIL DO
        WITH desc = SchemePair.Pair(p.first) DO
          res[Int(desc.first)] := Int(desc.rest)
        END;
        p := SchemePair.Pair(p.rest)
      END
    END;
    RETURN res
  END ToModula_A;
\end{lstlisting}

\textbf{Multidimensional arrays} are nested.  Modula-3's
\texttt{ARRAY [0..2], [0..3] OF REAL} is equivalent to
\texttt{ARRAY [0..2] OF ARRAY [0..3] OF REAL}---an array of arrays.
\texttt{sstubgen} generates converters recursively: the outer array's
element converter is itself the converter for the inner array.  In
Scheme, this becomes a list of lists:
\texttt{((0 . ((0 . 1.0) (1 . 2.0) ...)) (1 . ((0 . 3.0) ...)) ...)}.

\textbf{Open arrays} (\texttt{ARRAY OF T}, with no fixed bounds)
cannot exist as bare value types---they exist only as
\texttt{REF ARRAY OF T} or as formal parameters.  As \texttt{REF}
types, they are already \texttt{REFANY} and can be passed through the
Scheme boundary directly; \texttt{sstubgen} generates converters that
treat them as Scheme vectors or lists depending on context.

The key design choice is that \texttt{sstubgen} \emph{decomposes}
value types into Scheme-native structures rather than wrapping them in
opaque references.  This means Scheme code can inspect and construct
M3 records and arrays using ordinary Scheme operations (\texttt{car},
\texttt{cdr}, \texttt{assq}), at the cost of a copy at each boundary
crossing.


\subsection{How M3 Values Enter Scheme}
\label{sec:values-enter}

Because \texttt{SchemeObject.T = REFANY}, any M3 reference value can
be passed into Scheme simply by assigning it to a Scheme variable.
For a host program embedding MScheme:

\begin{lstlisting}[language=Modula3]
(* From the host M3 program: *)
VAR interp := NEW(SchemeM3.T).init(paths^);
interp.bind("my-driver", myM3Object);
\end{lstlisting}

The \texttt{bind} method defines a symbol in the global Scheme
environment.  The M3 object \texttt{myM3Object} becomes directly
accessible in Scheme as the value of \texttt{my-driver}.  No
conversion, copying, or indirection.  Scheme code manipulating this
value is manipulating the exact same M3 heap object.


\subsection{How Scheme Calls M3 Methods}
\label{sec:call-methods}

MScheme uses a three-layer dispatch mechanism:

\textbf{Layer~1: Type operations registry}
(\texttt{SchemeProcedureStubs.m3}, a hand-written module that provides
the registration framework; the individual dispatch procedures and
call stubs plugged into it are generated by \texttt{sstubgen}).  For
each M3 type that should have methods callable from Scheme, operations
are registered by typecode:

\begin{lstlisting}[language=Modula3]
SchemeProcedureStubs.RegisterOp(
    TYPECODE(TextSeq.T),    (* M3 runtime typecode *)
    "call-method",           (* operation name *)
    TextSeqCallMethod);      (* M3 procedure to dispatch *)
\end{lstlisting}

The dispatch procedure receives the interpreter, the M3 object, a
method name (as an Atom), and the Scheme argument list.  The
\texttt{interp} parameter is part of the standard signature for all
dispatch procedures and call stubs; it is passed through to the
individual method stubs, which may need it for exception handling or
Scheme callbacks, even when the dispatcher itself does not use it
directly:

\begin{lstlisting}[language=Modula3]
PROCEDURE TextSeqCallMethod(interp : Scheme.T;
                            object : SchemeObject.T;
                            args   : SchemeObject.T) : SchemeObject.T
  RAISES { Scheme.E } =
VAR name := Scheme.SymbolCheck(SchemeUtils.First(args));
    rest := SchemeUtils.Rest(args);
BEGIN
  IF name = addhi THEN
    NARROW(object, TextSeq.T).addhi(SchemeUtils.Str(SchemeUtils.First(rest)));
    RETURN SchemeBoolean.True()
  ELSIF name = size THEN
    RETURN SchemeLongReal.FromI(NARROW(object, TextSeq.T).size())
  ELSIF name = get THEN
    RETURN NARROW(object, TextSeq.T).get(SchemeLongReal.Int(SchemeUtils.First(rest)))
  (* ... *)
  END
END TextSeqCallMethod;
\end{lstlisting}

\textbf{Layer~2: Scheme-level wrapper} (\texttt{scheme-lib/src/m3.scm}).
The \texttt{obj-method-wrap} function creates a closure that dispatches
method calls through the type operation registry:

\begin{lstlisting}[language=Scheme]
(define (obj-method-wrap obj type)
  (lambda args
    (if (eq? (car args) '*m3*)
        obj    ; return the raw M3 object
        (modula-type-op type 'call-method obj (car args) (cdr args)))))
\end{lstlisting}

\textbf{Layer~3: Application code.}  The result is natural Scheme syntax:

\begin{lstlisting}[language=Scheme]
(define seq (obj-method-wrap (new-modula-object 'TextSeq.T) 'TextSeq.T))
(seq 'addhi "hello")
(seq 'size)           ;=> 1
(seq 'get 0)          ;=> "hello"
(seq '*m3*)           ;=> the raw TextSeq.T for passing back to M3
\end{lstlisting}

The \texttt{'*m3*} convention lets Scheme code extract the underlying
M3 object to pass it back into M3 methods that expect typed arguments.


\subsection{How Scheme Values Return to M3}
\label{sec:values-return}

Because Scheme values \emph{are} M3 \texttt{REFANY}, return is trivial.
An M3 method receiving a Scheme value simply narrows it:

\begin{lstlisting}[language=Modula3]
VAR result := interp.evalInGlobalEnv(someExpr);
(* result is REFANY -- narrow to expected type: *)
VAR text := NARROW(result, TEXT);
\end{lstlisting}

If the Scheme code returned a wrapped M3 object, the narrow succeeds
and the caller gets back the original M3 object with its original
type.  If it returned a Scheme-specific type (pair, number), the
narrow can distinguish that.  \texttt{ISTYPE} and \texttt{TYPECASE}
handle all cases safely.


\subsection{The Stub Generator (sstubgen)}
\label{sec:sstubgen}

MScheme includes a \emph{stub generator} that processes M3 interface
files at build time and auto-generates bidirectional binding code.
This is the system's most distinctive feature.

\texttt{sstubgen} shares substantial code with the Network Objects
stub generator (\texttt{m3-comm/stubgen})~\cite{birrell1994}.  Both
are \texttt{M3ToolFrame} extensions that walk the M3 AST using the M3
Toolkit (m3tk), extract type information via an \texttt{AstToType}
module, and generate code from the type signatures.  The network
objects \texttt{stubgen} generates marshalling and unmarshalling
procedures for remote method invocation; \texttt{sstubgen} adds
\texttt{TypeTranslator} and \texttt{ValueTranslator} modules that
generate \texttt{ToScheme}/\texttt{ToModula} converters and
\texttt{CallStub} wrappers for Scheme interoperability.  (See
Section~\ref{sec:lineage} for the provenance of \texttt{sstubgen} and
m3tk.)

The design goal, as stated in the \texttt{sstubgen} documentation, is:

\begin{quote}
``We want to be able to \emph{write Modula-3 code in Scheme.}  That
is, anything that can be coded in Modula-3 one should be able to code
in Scheme.''
\end{quote}

The stub generator operates in three phases:

\textbf{Phase~1} parses each M3 interface file using the M3 Toolkit
(m3tk) AST infrastructure.  For each procedure, type, constant, and
variable in the interface, it extracts type signatures and generates
an M3 stub module containing:

\begin{itemize}
\item \emph{ToModula} converters: functions that take a
  \texttt{SchemeObject.T} and produce the appropriate M3 type (with
  type-checking and bounds validation).
\item \emph{ToScheme} converters: functions that take an M3 value and
  produce a \texttt{SchemeObject.T}.
\item \emph{CallStub} procedures: wrappers that unpack Scheme
  arguments, call the real M3 procedure, convert the result, and
  handle exceptions.
\item Registration code that plugs these stubs into the runtime
  registry.
\end{itemize}

The generated code has two kinds of wrappers: \emph{call stubs} for
interface procedures and \emph{dispatch procedures} (with associated
\emph{method stubs}) for object method calls.

A \textbf{call stub} wraps a single M3 procedure.  It receives an
exception handler (\texttt{excH}) as an explicit parameter because the
caller (the Scheme evaluator) provides it directly.  The
\texttt{interp} parameter is needed to invoke the exception handler
via \texttt{SchemeApply} if the M3 procedure raises.  A generated call
stub looks like:

\begin{lstlisting}[language=Modula3]
PROCEDURE CallStub_Text_Cat(interp : Scheme.T;
                            args   : SchemeObject.T;
                            excH   : SchemeObject.T) : SchemeObject.T
  RAISES { Scheme.E } =
VAR p := SchemePair.Pair(args);
    a1 := ToModula_TEXT(First(p));   (* Scheme -> TEXT *)
    a2 := ToModula_TEXT(Second(p));
BEGIN
  RETURN ToScheme_TEXT(Text.Cat(a1, a2))  (* TEXT -> Scheme *)
END CallStub_Text_Cat;
\end{lstlisting}

A \textbf{dispatch procedure} routes method calls by name.  It does
not take \texttt{excH} as a separate parameter; instead, the exception
handler is packed into the \texttt{args} list (as the third element)
by the Scheme-level \texttt{obj-method-wrap} caller, and the
dispatcher extracts it before forwarding to the appropriate method
stub.  The method stubs themselves have the same signature as call
stubs, including the \texttt{excH} parameter.

\textbf{Phase~2} collects all stub modules within a package into a
single \texttt{RegisterStubs()} entry point.

\textbf{Phase~3} links all package-level \texttt{RegisterStubs} into
the final executable, so that a single call at startup activates all
bindings.

The m3makefile for the CSP compiler (cspc) shows this in practice:

\begin{lstlisting}[language={}]
import ("sstubgen")
SchemeStubs ("CspExpression")
SchemeStubs ("Fmt")
SchemeStubs ("Atom")
SchemeStubs ("Word")
SchemeStubs ("M3Ident")
SchemeStubs ("FileWr")
(* ... 30+ more interfaces ... *)
ExportSchemeStubs ("cspc")
\end{lstlisting}

Each \texttt{SchemeStubs} invocation processes one M3 interface.  The
result is that Scheme code in cspc can call \texttt{Fmt.Int},
\texttt{Atom.FromText}, \texttt{FileWr.Open}, etc., as naturally as
M3 code can.


\subsection{Call Flow for Interface Procedures}
\label{sec:call-flow}

When Scheme code evaluates \texttt{(Wr.PutText wr "hello")}, the call
passes through four layers before reaching the actual M3 procedure.
Tracing the full path:

\textbf{1.\ Evaluation.}  \texttt{Scheme.Eval} in \texttt{Scheme.m3}
evaluates the operator position.  The symbol \texttt{Wr.PutText}
resolves to a \texttt{SchemePrimitive.T} object that was registered at
startup.  The two argument expressions are evaluated, producing
\texttt{SchemeObject.T} values.  For two arguments, the evaluator
calls the optimized \texttt{apply2} method directly.

\textbf{2.\ Primitive dispatch.}
\texttt{SchemePrimitive.Apply2} validates the argument count and calls
the main dispatcher (\texttt{Prims}).  Interface procedure stubs are
registered with IDs above the range of the built-in primitive enum
\texttt{P}, so the dispatcher recognizes them as extensions and calls
through to the \texttt{ExtDefiner}:

\begin{lstlisting}[language=Modula3]
IF t.id > ORD(LAST(P)) THEN
  RETURN NARROW(t.definer, ExtDefiner).apply(t, interp, args)
\end{lstlisting}

\textbf{3.\ Stub lookup.}  \texttt{ExtDefiner.apply} looks up the
procedure ID in a hash table, finds a \texttt{PrimRec} containing the
qualified name (\texttt{Wr}, \texttt{PutText} as atoms), and calls
\texttt{SchemeProcedureStubs.CallStubApply}.  This function walks a
linked list of registered stubs, matching the interface name and item
name by atom identity (pointer comparison), and calls the matching
stub procedure.

\textbf{4.\ Stub execution.}  The call stub---generated by
\texttt{sstubgen} at build time---unpacks arguments from the Scheme
pair list, converts each from \texttt{SchemeObject.T} to the
appropriate M3 type using the \texttt{ToModula} functions, calls the
real M3 procedure, converts the result back via \texttt{ToScheme}, and
handles any M3 exceptions.  The \texttt{interp} parameter is needed to
invoke the Scheme-level exception handler via \texttt{SchemeApply} if
the M3 procedure raises; the \texttt{excH} parameter is the exception
handler closure itself:

\begin{lstlisting}[language=Modula3]
PROCEDURE CallStub_Wr_PutText(interp: Scheme.T;
    args, excH: SchemeObject.T): SchemeObject.T RAISES {Scheme.E} =
  VAR wr  := NARROW(First(args), Wr.T);           (* REFANY -> Wr.T *)
      txt := SchemeString.ToText(Second(args));    (* SchemeString -> TEXT *)
  BEGIN
    TRY Wr.PutText(wr, txt);                      (* actual M3 call *)
        RETURN SchemeBoolean.True()
    EXCEPT Wr.Failure(e) =>
        EVAL SchemeApply.OneArg(interp, excH, ...)
    END
  END CallStub_Wr_PutText;
\end{lstlisting}

(All generated call stubs follow the
\texttt{CallStub\_\textit{Interface}\_\textit{Procedure}} naming
convention; method stubs use
\texttt{MethodStub\_\textit{Type}\_\textit{Method}}.)

For the \texttt{wr} argument, which is already a traced reference
(\texttt{Wr.T} is an object type), the conversion is just a
\texttt{NARROW}---a type-checked pointer cast, no data copying.  For
the \texttt{txt} argument, \texttt{ToText} converts the mutable
Scheme string (\texttt{REF ARRAY OF CHAR}) to an immutable M3
\texttt{TEXT}.  The return value is a Scheme boolean.

The entire path---from \texttt{Eval} to the actual \texttt{Wr.PutText}
call---is native M3 procedure calls.  There is no interpreter loop,
bytecode dispatch, or reflection.  The overhead beyond a direct M3
call is: argument list traversal (walking \texttt{SchemePair} cells),
type conversion (one \texttt{NARROW} + one \texttt{ToText}), and stub
lookup (atom comparison in a linked list, typically short).


\subsection{Runtime Type Introspection from Scheme}
\label{sec:introspection}

MScheme exposes M3's runtime type system to Scheme via primitives:

\begin{lstlisting}[language=Scheme]
(rttype-typecode obj)          ; get an object's runtime typecode
(rttype-supertype tc)          ; get a typecode's supertype
(rttype-issubtype? tc1 tc2)    ; subtype relationship
(rttype-maxtypecode)           ; highest allocated typecode
(rtbrand-getname tc)           ; get the brand string for a type
(modula-type-class tc)         ; classification (object, ref, etc.)
(list-modula-type-ops tc)      ; list registered operations
(list-modula-types)            ; list all registered type names
(typename->typecode 'Foo.T)    ; look up typecode by name
\end{lstlisting}

This allows Scheme code to introspect M3 objects, walk the type
hierarchy, and discover available methods---something that M3 code
itself can do only in limited ways.


\subsection{Concurrency Model}
\label{sec:concurrency}

MScheme is designed for multi-threaded M3 programs.  The key design
rule: a single \texttt{Scheme.T} interpreter instance is
single-threaded, but multiple interpreters may share a global
environment.  A global environment is a
\texttt{SchemeEnvironment.T}---the root of a chain of lexical scoping
frames that holds all top-level bindings (primitives, stub
registrations, and user-defined globals).  Environments come in
\texttt{Safe} (synchronized) and \texttt{Unsafe} (unsynchronized)
variants; sharing a global environment across threads requires the
\texttt{Safe} variant.


\subsection{C Libraries via M3 as Glue}
\label{sec:c-libraries}

MScheme can be used even when the goal is simply to call C code from
Scheme.  Modula-3 serves as a type-safe glue layer: a thin M3
interface wraps the C functions, and \texttt{sstubgen} auto-exports
that interface to Scheme.  The programmer never writes
Scheme-specific binding code.

A concrete example is the \texttt{Mpz} module in the \texttt{m3utils}
tree, which wraps the GNU Multiple Precision Arithmetic Library (GMP)
for use from the CSP compiler's Scheme environment.  The architecture
has three layers:

\textbf{Layer~1: C binding.}  \texttt{MpzP.i3} declares GMP functions
as \texttt{EXTERNAL} procedures with C calling convention:

\begin{lstlisting}[language=Modula3]
(* MpzP.i3 -- raw C interface *)
<*EXTERNAL "__gmpz_add"*>
PROCEDURE c_add(f0: MpzPtrT; f1: MpzPtrT; f2: MpzPtrT);

<*EXTERNAL "__gmpz_mul"*>
PROCEDURE c_mul(f0: MpzPtrT; f1: MpzPtrT; f2: MpzPtrT);

<*EXTERNAL "mpz_format_decimal"*>
PROCEDURE format_decimal(z: MpzPtrT): Ctypes.const_char_star;
\end{lstlisting}

A small C file (\texttt{MpzC.c}) provides helper functions for
formatting via \texttt{gmp\_asprintf}.

\textbf{Layer~2: Type-safe M3 wrapper.}  \texttt{Mpz.i3} defines an
opaque traced type \texttt{T <: REFANY} that wraps the raw
\texttt{mpz\_t}.  The implementation (\texttt{MpzOps.m3}) extracts the
address and calls through to layer~1:

\begin{lstlisting}[language=Modula3]
(* Mpz.i3 -- public interface *)
TYPE T <: REFANY;
PROCEDURE New(): T;
PROCEDURE add(f0: T; f1: T; f2: T);
PROCEDURE Format(t: T; base := FormatBase.Decimal): TEXT;
\end{lstlisting}

The concrete type is revealed in \texttt{MpzRep.i3}:

\begin{lstlisting}[language=Modula3]
(* MpzRep.i3 -- concrete revelation *)
REVEAL
  Mpz.T = BRANDED Mpz.Brand OBJECT
    val : Val;
  END;

TYPE
  Val = RECORD
    w : ARRAY [0..2] OF Word.T;
    (* this will be cast to an mpz_t, and it is AT LEAST big enough. *)
  END;
\end{lstlisting}

The \texttt{val} field embeds the GMP \texttt{mpz\_t} storage directly
in the M3 heap object.  The implementation extracts its address for C
calls:

\begin{lstlisting}[language=Modula3]
(* MpzOps.m3 -- implementation *)
PROCEDURE add(f0: T; f1: T; f2: T) =
  BEGIN P.c_add(ADR(f0.val), ADR(f1.val), ADR(f2.val)) END add;
\end{lstlisting}

Because \texttt{T <: REFANY}, \texttt{Mpz.T} values are traced by the
M3 garbage collector.  The opaque type prevents Scheme (or M3) code
from accessing the raw \texttt{mpz\_t} memory directly---only the
operations declared in \texttt{Mpz.i3} are available.

Resource cleanup uses Modula-3's \texttt{WeakRef} facility.  When
\texttt{Mpz.New} creates a new wrapper object, it registers a
weak-reference callback via \texttt{WeakRef.FromRef}:

\begin{lstlisting}[language=Modula3]
(* Mpz.m3 *)
PROCEDURE New() : T =
  VAR res := NEW(T);
  BEGIN
    P.c_init(LOOPHOLE(ADR(res.val), P.MpzPtrT));
    EVAL WeakRef.FromRef(res, CleanUp);
    set_ui(res, 0);
    RETURN res
  END New;

PROCEDURE CleanUp(<*UNUSED*>READONLY w : WeakRef.T; r : REFANY) =
  BEGIN
    WITH this = NARROW(r, T) DO
      P.c_clear(ADR(this.val))
    END
  END CleanUp;
\end{lstlisting}

When the GC determines that an \texttt{Mpz.T} object is unreachable,
it schedules the \texttt{CleanUp} callback, which calls GMP's
\texttt{mpz\_clear} to release the C-level memory.  No explicit
deallocation is needed from either M3 or Scheme code---the
\texttt{Mpz.T} values can be passed around, stored in lists, and
forgotten, just like any other Scheme value.  This is the same
\texttt{WeakRef} mechanism that Modula-3 programs use for any C
resource: FFTW plans, file descriptors, database handles.

\textbf{Layer~3: Automatic Scheme export.}  Adding
\texttt{SchemeStubs("Mpz")} to the m3makefile causes \texttt{sstubgen}
to generate Scheme bindings for every procedure in \texttt{Mpz.i3}.
From Scheme:

\begin{lstlisting}[language=Scheme]
(define a (Mpz.New))
(define b (Mpz.NewInt 42))
(define c (Mpz.New))
(Mpz.add c a b)
(Mpz.Format c)    ;=> "42"
\end{lstlisting}

The \texttt{Mpz.T} values flowing through Scheme are the same traced
M3 heap objects---no copying or marshaling.  The GC handles their
lifetime.  The C library's \texttt{mpz\_t} storage is embedded inside
the M3 object's record fields and is deallocated when the M3 object
is collected.

This pattern---C library, M3 interface, \texttt{sstubgen}---is
general.  The same approach could wrap SQLite, zlib, OpenSSL, or any
C library.  Modula-3's \texttt{EXTERNAL} pragma provides the C
binding; its type system provides safety; \texttt{sstubgen} provides
the Scheme bridge.  The result is that MScheme functions as a
practical Scheme-with-C-FFI system, even though its design is not
about C interoperability at all.

The \texttt{Mpz} module is itself partly auto-generated: a Scheme
script (\texttt{make-mpz.scm}) reads GMP's function prototypes and
emits \texttt{MpzP.i3}, \texttt{Mpz.i3}, and
\texttt{MpzOps.m3}---over 200 functions.  This is a second level of
code generation: Scheme generates M3 code, then \texttt{sstubgen}
generates Scheme stubs from that M3 code, closing the loop.


\subsection{Summary of Key Properties}
\label{sec:summary}

\begin{itemize}
\item \textbf{Zero-copy value sharing}: M3 objects are Scheme objects,
  no wrapping needed.
\item \textbf{Unified GC}: Both languages share the M3 traced heap
  and garbage collector.  No reference counting, prevent-collection
  APIs, or weak reference coordination.
\item \textbf{Automatic stub generation}: Build-time processing of M3
  interfaces produces bidirectional bindings with type checking.
\item \textbf{Runtime type introspection}: Scheme code can query the
  M3 type system, discovering types and operations dynamically.
\item \textbf{Method dispatch via typecode}: $O(n)$ scan of
  registered operations per typecode, then atom comparison for method
  name.
\end{itemize}


%% ==================================================================
\section{Comparable Systems}
\label{sec:comparable}
%% ==================================================================

\subsection{Lua/C (Lua C API)}
\label{sec:lua}

\textbf{Core representation}: A \emph{stack} of tagged union values
(\texttt{TValue} internally).  The C API never exposes \texttt{TValue}
directly; all value exchange is through positional stack slots.

\textbf{Wrapping C objects}: \texttt{lua\_newuserdata(L, size)}
allocates a GC-managed memory block.  A \emph{metatable} (a Lua
table) is attached to define behavior.  The \texttt{\_\_index}
metamethod routes field/method lookups to a table of C functions.
\texttt{\_\_gc} provides a destructor.

\textbf{Method dispatch}: When Lua code calls
\texttt{obj:method(args)}, the VM performs a metatable lookup for
\texttt{\_\_index}, finds the method in the method table, and calls
the C function with the Lua state.  The C function pops arguments from
the stack and pushes results.

\textbf{Values back to C}: Typed accessors read stack slots:
\texttt{lua\_tonumber}, \texttt{lua\_tostring},
\texttt{lua\_touserdata}.

\textbf{Registration}: C functions are registered via arrays of
\texttt{luaL\_Reg} structs, or individually with
\texttt{lua\_pushcfunction}.

\textbf{Key trade-offs}: The stack-based API provides complete ABI
isolation---Lua's internal value representation can change without
breaking extensions.  But every boundary crossing requires explicit
push/pop operations.  There is no automatic type mapping, stub
generation, or introspection of the host's type system.


\subsection{GNU Guile (Scheme in C)}
\label{sec:guile}

\textbf{Core representation}: \texttt{SCM}, a machine-word-sized
tagged value.  Fixnums, characters, and booleans are encoded in tag
bits.  Heap objects (pairs, foreign objects, strings) are tagged
pointers.

\textbf{Wrapping C objects}: \emph{Foreign object types} replace the
older SMOB mechanism.
\texttt{scm\_make\_foreign\_object\_type} defines a type with named
slots and a finalizer.  Instances store C pointers in slots.

\textbf{Method dispatch}: Typically via named procedures
(\texttt{my-object-get-x}).  Optional GOOPS integration (CLOS-style
generic functions with multiple dispatch) is possible but uncommon for
foreign objects.

\textbf{Values back to C}: \texttt{scm\_to\_int(val)},
\texttt{scm\_to\_double(val)}, etc.  Every function takes and returns
\texttt{SCM}.

\textbf{Registration}:
\texttt{scm\_c\_define\_gfunc("name", nreq, nopt, nrest, cfunc)}
registers a C function as a Scheme procedure in the current module.

\textbf{Key trade-offs}: The uniform \texttt{SCM} type is simpler
than Lua's stack (no positional bookkeeping), but every C function
must do its own type checking.  Guile's conservative GC (BDW-GC) is
very friendly to C embedders---no need to track roots manually.  But
Guile provides no automatic stub generation from C header files.


\subsection{Tcl/C}
\label{sec:tcl}

\textbf{Core representation}: \texttt{Tcl\_Obj*}, a dual-ported value
with both a string representation and a cached native (``internal'')
representation.  Conversions between the two forms are lazy and
cached.  Reference counted (no GC).

\textbf{Wrapping C objects}: Tcl uses the ``command-as-object''
pattern.  Creating an object registers a new Tcl command named after
it.  The \texttt{ClientData} (a \texttt{void*}) stored with the
command points to the C struct.  ``Methods'' are subcommands:
\texttt{.button configure -text "OK"}.

\textbf{Method dispatch}: String-based.  The first argument to the
object's command is the subcommand name, matched by string comparison
or hash lookup.  Object identity is a string name.

\textbf{Registration}:
\texttt{Tcl\_CreateObjCommand(interp, "name", func, clientData,
  deleteProc)}.

\textbf{Key trade-offs}: ``Everything is a string'' gives
extraordinary composability but loses type safety.  Reference counting
provides deterministic destruction but cannot handle cycles.  The
command-as-object pattern is elegant and uniform but means object
identity is a string, and there is no direct pointer from Tcl to C
memory without going through command lookup.


\subsection{CPython Extension API}
\label{sec:cpython}

\textbf{Core representation}: \texttt{PyObject*}, a heap-allocated
struct with a reference count and a pointer to a
\texttt{PyTypeObject}.

\textbf{Wrapping C objects}: The embedder fills out a
\texttt{PyTypeObject} struct with dozens of slots:
\texttt{tp\_methods} (method table), \texttt{tp\_getset} (computed
properties), \texttt{tp\_members} (direct struct field access by
offset), \texttt{tp\_dealloc} (destructor), plus slots for numeric,
sequence, mapping, iterator, and async protocols.

\textbf{Method dispatch}: Attribute lookup via the type's method
resolution order (MRO).  If the name is found in
\texttt{tp\_methods}, a bound method object is created wrapping the C
function.

\textbf{Values back to C}: \texttt{PyLong\_AsLong(obj)},
\texttt{PyFloat\_AsDouble(obj)},
\texttt{PyArg\_ParseTuple(args, "is|d", ...)}.

\textbf{Registration}: Module initialization fills a
\texttt{PyModuleDef} struct and registers types with
\texttt{PyType\_Ready}.

\textbf{Key trade-offs}: Maximum control over Python's internals, but
the reference counting discipline is notoriously error-prone.  Every
boundary crossing allocates heap objects.  Defining even a simple type
requires filling many \texttt{PyTypeObject} slots.  The API is
primarily designed for \emph{extending} Python, not for
\emph{embedding} it.


\subsection{Kawa Scheme (JVM)}
\label{sec:kawa}

\textbf{Core representation}: \texttt{java.lang.Object}.  Kawa
compiles Scheme to JVM bytecode; Scheme values are Java objects.

\textbf{Wrapping Java objects}: No wrapping needed.  Java objects are
directly usable in Kawa Scheme.
\texttt{(define pt (java.awt.Point 10 20))} creates a Java
\texttt{Point}.  \texttt{pt:x} accesses the field.

\textbf{Method dispatch}: When the type is known at compile time, Kawa
generates \texttt{invokevirtual} bytecode---as efficient as Java.
When unknown, it falls back to reflection.  Type annotations
(\texttt{::java.awt.Point}) guide the compiler.

\textbf{Values back to Java}: Scheme values are already JVM objects.
Java code invokes Scheme functions by calling \texttt{apply} on
\texttt{gnu.mapping.Procedure} objects.

\textbf{Registration}: None needed---both languages share the JVM
class loader and object space.

\textbf{Key trade-offs}: Zero-cost interoperability when types are
known, but reflection fallback is slow.  Kawa must work within JVM
constraints (no true tail calls or first-class continuations without
trampolining).  The shared-runtime model is fundamentally different
from all other systems here.


\subsection{JScheme (Scheme in Java)}
\label{sec:jscheme}

\textbf{Core representation}: \texttt{java.lang.Object}.  JScheme is
a tree-walking Scheme interpreter written in Java.  In Norvig's
original (1998)~\cite{norvig1998}, all numbers are
\texttt{java.lang.Double}, symbols are interned
\texttt{java.lang.String}, strings are \texttt{char[]} (for
mutability), pairs are a custom \texttt{Pair} class with
\texttt{Object first, rest} fields, and \texttt{null} represents the
empty list.  The later Anderson/Hickey version
(2003+)~\cite{anderson2003} switched numbers to \texttt{Integer} or
\texttt{Double} as appropriate, symbols to a custom \texttt{Symbol}
class, and strings to \texttt{java.lang.String}.

\textbf{Wrapping Java objects}: No wrapping needed.  Any Java object
is already an \texttt{Object} and therefore a valid Scheme value.  The
evaluator treats anything that is not a symbol and not a pair as
self-evaluating, so a \texttt{java.awt.Frame} or
\texttt{java.util.HashMap} placed into a Scheme variable simply
evaluates to itself.

\textbf{Method dispatch}: Entirely reflection-based.  The later
JScheme provides a dot-notation syntactic sugar parsed at read time:

\begin{lstlisting}[language=Scheme]
(import "java.awt.Font")
(Font. "Serif" Font.BOLD$ 24)     ; constructor (trailing dot)
(.setFont component font)         ; instance method (leading dot)
(.name$ obj)                      ; field access (trailing $)
(Math.round 123.456)              ; static method (dot in middle)
\end{lstlisting}

These desugar into calls to \texttt{JavaMethod},
\texttt{JavaConstructor}, or \texttt{JavaField} objects, all of which
use \texttt{java.lang.reflect} internally.  Method tables are cached
per class in hash tables, and overload resolution follows Java's
``most specific applicable method'' rule at runtime.

\textbf{Values back to Java}: Scheme values are already
\texttt{Object}.  Java code calls
\texttt{JScheme.call("proc-name", arg1, arg2)} and receives an
\texttt{Object} result, narrowing via cast.

\textbf{Scheme as Java interface}: The
\texttt{SchemeInvocationHandler} class uses
\texttt{java.lang.reflect.Proxy} to let a Scheme closure implement
any Java interface, so Scheme callbacks work anywhere Java expects an
interface (listeners, \texttt{Runnable}, etc.).

\textbf{Registration}: \texttt{(import "java.awt.*")} makes classes
available; dot-notation symbols are resolved against the import table.
From the Java side, \texttt{JScheme.setGlobalValue("name", obj)}
injects values.

\textbf{Key trade-offs}: JScheme achieves the same zero-wrapping
interoperability as Kawa---Java objects \emph{are} Scheme values---but
pays for it with reflection overhead on every method call.  Kawa
compiles to \texttt{invokevirtual} bytecode; JScheme interprets and
reflects.  The simplicity is striking (Norvig's original is
\textasciitilde1700 lines), but the performance gap is large.  JScheme
is the direct ancestor of MScheme and establishes the key principle
that MScheme inherits: when the host language has a universal reference
type and a GC, the embedding boundary can be made nearly invisible.


\subsection{Common Lisp / CFFI}
\label{sec:cffi}

\textbf{Core representation}: Tagged Lisp objects
(implementation-dependent tagging scheme).  Foreign values are
\texttt{cffi:foreign-pointer} wrappers.

\textbf{Wrapping C objects}: The programmer wraps foreign pointers in
CLOS objects and defines generic functions that call through to C.
\texttt{trivial-garbage:finalize} attaches destructors.

\textbf{Method dispatch}: Generic functions with CLOS multiple
dispatch.  Methods dispatch on the classes of \emph{all} arguments,
not just the receiver.

\textbf{Values back to C}:
\texttt{cffi:foreign-slot-value} reads struct fields by declared
offset.  \texttt{defcfun} defines C function bindings with automatic
marshaling.

\textbf{Registration}: \emph{Pull-based}---the Lisp side declares C
bindings with \texttt{defcfun}, \texttt{defcstruct},
\texttt{defcenum}.  The C library is unmodified.

\textbf{Key trade-offs}: CFFI is portable across CL implementations
but requires manual struct layout declarations and memory management.
CLOS provides a powerful ``glue'' layer but is per-project, not
standardized for FFI.


\subsection{Racket FFI}
\label{sec:racket-ffi}

\textbf{Core representation}: Tagged Racket objects.  C types are
first-class Racket values: \texttt{\_int}, \texttt{\_double},
\texttt{\_pointer}, \texttt{\_fun}.

\textbf{Wrapping C objects}:
\texttt{define-cpointer-type} creates tagged pointer types.
\texttt{define-cstruct} auto-generates constructors, accessors, and
mutators.  Custom type combinators (\texttt{make-ctype}) attach
conversion functions to base C types.

\textbf{Method dispatch}: Not object-oriented at the FFI layer.
Operations on foreign values are called as procedures.

\textbf{Values back to C}: \texttt{\_fun} type descriptors handle all
marshaling declaratively.
\texttt{(\_fun \_string \_int -> \_bool)} specifies argument and
return conversions.

\textbf{Registration}: \emph{Pull-based}---\texttt{get-ffi-obj} or
\texttt{define-ffi-definer} declares bindings.  The C library is
unmodified.

\textbf{Key trade-offs}: The most algebraic and composable FFI design.
Types as values enable declarative marshaling.  Tagged pointers add
safety.  But the C library must be described entirely from the Racket
side, and per-call marshaling overhead is significant for tight loops.


\subsection{Tinylisp (Modula-2+)}
\label{sec:tinylisp}

\textbf{Core representation}: \texttt{REFANY}.  Tinylisp (DEC Systems
Research Center, 1989)~\cite{ellis1989} is a Lisp implemented in
Modula-2+, the SRC dialect of Modula-2 that added garbage collection,
exception handling, \texttt{REFANY}, and \texttt{TYPECASE}.  Symbolic
expressions are ordinary Modula-2+ ref types: \texttt{Ref.Integer},
\texttt{Ref.LongReal}, \texttt{Ref.Boolean}, \texttt{Ref.Char},
\texttt{Text.T}, \texttt{SxSymbol.T}, \texttt{List.T},
\texttt{Ref.Vector}.  Any application's opaque-ref type is directly
usable as a Tinylisp value---a \texttt{Wire.T} or \texttt{Wr.T} is a
first-class s-expression, just as an \texttt{Mpz.T} is a first-class
MScheme value.

\textbf{Wrapping Modula-2+ objects}: No wrapping needed.  Any traced
reference is already \texttt{REFANY} and therefore a valid Tinylisp
value.

\textbf{Method dispatch}: Via generated stubs.  A \emph{stub compiler}
(\texttt{compile\_tli}) processes hand-written \texttt{.tli} (TinyLisp
Interface) files and generates Modula-2+ stub procedures that
dynamically check argument counts and types, convert between
s-expression types and Modula-2+ value types, and call the
corresponding procedures.  A \texttt{.tli} file declares modules,
types, exceptions, and procedures in s-expression syntax:

\begin{lstlisting}[language=Scheme]
(MODULE Text Lisp)
(TYPE T)
(EXCEPTION EndOfFile)
(PROCEDURE FromChar (CHAR) "Text.T")
(PROCEDURE Length ("Text.T") INTEGER)
(PROCEDURE SubText ("Text.T" INTEGER & INTEGER) "Text.T")
\end{lstlisting}

\textbf{Values back to Modula-2+}: Scheme values are already
\texttt{REFANY}.  The stubs narrow them via \texttt{TYPECASE} for
value types (\texttt{INTEGER}, \texttt{CHAR}, \texttt{BOOLEAN}, etc.)
and \texttt{NARROW} for reference types.

\textbf{Registration}: The generated stubs register procedures at
startup, similar to MScheme's approach.

\textbf{Execution model}: Unlike MScheme's tree-walking interpreter,
Tinylisp compiled expressions directly into machine code at runtime
(an on-the-fly compiler), making Tinylisp-to-Tinylisp calls roughly
as fast as Modula-2+ calls.

\textbf{Key trade-offs}: Tinylisp achieves the same zero-wrapping,
shared-heap architecture as MScheme---the underlying mechanism
(\texttt{REFANY} + tracing GC + \texttt{TYPECASE}) is identical.  The
key difference is in stub generation: Tinylisp requires a hand-written
\texttt{.tli} file (essentially a manual transliteration of the
Modula-2+ \texttt{.def} file with type annotations), while MScheme's
\texttt{sstubgen} reads the standard M3 \texttt{.i3} interface files
directly via m3tk.  Tinylisp is also faster at the Lisp level
(compiled vs.\ interpreted) but was designed as a lightweight extension
language rather than a full Scheme.


%% ==================================================================
\section{Comparative Analysis}
\label{sec:analysis}
%% ==================================================================

\subsection{The Central Trade-off: Shared Heap vs.\ Marshaling Boundary}
\label{sec:shared-heap}

The most fundamental architectural difference across these systems is
whether the host and guest languages \emph{share} their object
representations or maintain \emph{separate} representations connected
by a marshaling boundary.

\textbf{Shared heap} (MScheme, Tinylisp, JScheme, Kawa): The host
language's objects are directly usable as scripting language values.
No per-call conversion cost.  The garbage collector sees everything.
This is only possible when both languages are GC'd and share a runtime
(M3 + MScheme, Modula-2+ + Tinylisp) or a VM (Java + JScheme, Java +
Kawa).

\textbf{Marshaling boundary} (Lua, Guile, Tcl, Python, CFFI, Racket):
The host creates scripting-language values via conversion APIs.  Every
boundary crossing allocates and converts.  The two sides maintain
separate type systems and (often) separate memory management.

MScheme's \texttt{SchemeObject.T = REFANY} and Tinylisp's use of
\texttt{REFANY} are the extreme cases of shared heap: it is not that
host objects are \emph{wrapped} as Lisp/Scheme values, but that they
literally \emph{are} Lisp/Scheme values, by type identity.


\subsection{Type Safety and Dispatch Cost}
\label{sec:type-safety}

\begin{table}[ht]
\centering
\small
\begin{tabular}{@{}lll@{}}
\toprule
System   & Type check at boundary                 & Method dispatch cost \\
\midrule
MScheme  & \texttt{TYPECASE} (M3 runtime typecheck)      & Atom comparison per registered op \\
Tinylisp & \texttt{TYPECASE} (M2+ runtime typecheck)     & Stub dispatch (compiled) \\
Lua      & Metatable lookup + manual C-side checks & Hash lookup in metatable \\
Guile    & Manual \texttt{SCM} type predicates           & Symbol lookup in module \\
Tcl      & String parse on demand                 & String comparison for subcommand \\
Python   & \texttt{PyArg\_ParseTuple} format string       & MRO attribute lookup \\
Kawa     & JVM \texttt{checkcast} instruction            & \texttt{invokevirtual} or reflection \\
JScheme  & \texttt{instanceof} (runtime)                 & Reflection + cached method table \\
CFFI     & CLOS \texttt{typep} + manual checks           & Generic function dispatch \\
Racket   & cpointer tag check                     & Direct call (no dispatch) \\
\bottomrule
\end{tabular}
\caption{Type safety and dispatch cost across systems.}
\label{tab:type-safety}
\end{table}

MScheme benefits from M3's runtime type system: \texttt{TYPECASE} is a
compiler-supported, type-safe downcast with $O(1)$ performance (the
runtime checks the typecode, which is an integer).  This is
essentially the same mechanism as Java's \texttt{instanceof} check,
used by Kawa and JScheme on the JVM.


\subsection{Automatic Binding Generation}
\label{sec:auto-binding}

\begin{table}[ht]
\centering
\small
\begin{tabular}{@{}llll@{}}
\toprule
System   & Auto-generation? & Input & Mechanism \\
\midrule
MScheme  & Yes (\texttt{sstubgen})       & M3 \texttt{.i3} interface files    & M3 Toolkit AST parsing + code gen \\
Tinylisp & Yes (\texttt{compile\_tli})   & Hand-written \texttt{.tli} files   & Stub compiler generates M2+ stubs \\
Lua      & No                     & ---                           & Manual C code \\
Guile    & Partial (g-wrap)       & C headers                   & Third-party tool, not standard \\
Tcl      & Partial (critcl, SWIG) & C headers                   & Third-party tools \\
Python   & Partial (Cython, etc.) & C/C++ headers + annotations & Third-party tools \\
Kawa     & Unnecessary            & ---                           & Shared JVM class loader \\
JScheme  & Unnecessary            & ---                           & Reflection on JVM classes at runtime \\
CFFI     & No                     & Manual declarations         & Programmer writes \texttt{defcfun} \\
Racket   & No                     & Manual declarations         & Programmer writes \texttt{define-ffi-definer} \\
\bottomrule
\end{tabular}
\caption{Automatic binding generation across systems.}
\label{tab:auto-binding}
\end{table}

MScheme's \texttt{sstubgen} is notable because it is a
\emph{first-party} tool that ships with the interpreter, processes the
host language's standard interface files, and generates complete
bidirectional bindings.  The equivalent in the Python world would be a
tool that reads \texttt{.pyi} stub files and generates complete C
extension modules automatically.

The closest comparable systems are SWIG (which processes C/C++ headers
to generate bindings for multiple target languages) and Kawa (which
needs no generator at all because Java and Scheme share the JVM).


\subsection{GC Coordination}
\label{sec:gc-coordination}

\begin{table}[ht]
\centering
\small
\begin{tabular}{@{}llll@{}}
\toprule
System   & Host GC            & Guest GC                   & Coordination \\
\midrule
MScheme  & M3 traced refs     & Same GC                    & None needed---unified heap \\
Tinylisp & M2+ traced refs    & Same GC                    & None needed---unified heap \\
Lua      & None (C)           & Incremental mark-sweep     & Explicit: stack + registry roots \\
Guile    & None (C)           & Conservative (BDW-GC)      & Automatic: scans C stack conservatively \\
Tcl      & None (C)           & Refcounting                & Manual: \texttt{Tcl\_IncrRefCount}/\texttt{Tcl\_DecrRefCount} \\
Python   & None (C)           & Refcounting + cycle detect & Manual: \texttt{Py\_INCREF}/\texttt{Py\_DECREF} \\
Kawa     & JVM GC             & Same GC                    & None needed---unified heap \\
JScheme  & JVM GC             & Same GC                    & None needed---unified heap \\
CFFI     & CL GC              & Same GC (for CL objects)   & \texttt{trivial-garbage:finalize} for foreign ptrs \\
Racket   & Precise moving GC  & Same GC                    & Special APIs for C extensions holding Racket refs \\
\bottomrule
\end{tabular}
\caption{GC coordination across systems.}
\label{tab:gc-coordination}
\end{table}

The MScheme/M3, Tinylisp/M2+, Kawa/JVM, and JScheme/JVM systems
share a critical advantage: since both languages use the same garbage
collector, there is no coordination problem.  A Scheme closure that
captures an M3 object keeps it alive automatically.  An M3 data
structure containing Scheme pairs is traced correctly.  No reference
counting, prevent-collection guards, or weak-reference bridges.

This eliminates an entire class of bugs that plagues C-based embedding
systems.  In CPython, getting reference counts wrong is the single
most common source of C extension bugs.  In Lua, forgetting to anchor
a value on the stack or in the registry causes it to be collected while
C still holds a pointer to it.

The fundamental problem runs deeper than API discipline, and
understanding it clarifies why the shared-heap architecture is confined
to a small set of languages: C makes garbage collection nearly
impossible to do correctly.  The language permits pointers to be hidden
in ways that are both legal and common in practice, defeating any
collector that tries to find them.

A ``conservative'' collector like BDW-GC works by scanning the C
stack, registers, and static data segments for bit patterns that look
like they might be heap addresses.  This is inherently a heuristic,
and it fails in both directions.

\textbf{False negatives (hidden pointers---the dangerous case).}  C
allows a program to store a pointer as an integer
(\texttt{uintptr\_t}), XOR two pointers together (XOR linked lists),
offset a pointer by a constant and recover it later, split a pointer's
bytes across non-adjacent struct fields, send it through a pipe or
socket and read it back, or store it in a memory-mapped file.  In all
of these cases, the pointer exists but does not appear on the stack or
in a register in a form the collector can recognize.  If the collector
runs at that moment, it may conclude the object is unreachable and
free it.  The program then recovers the hidden pointer and
dereferences freed memory.  This is not a contrived scenario---XOR-linked
lists save memory in embedded systems, pointer tagging is used in
language runtimes (including CPython's own \texttt{ob\_refcnt} field
packing), and integer casts are routine in serialization and IPC code.

\textbf{False positives (pointer lookalikes---the slow leak).}  On a
32-bit system, roughly 1 in 1000 random integers looks like a valid
heap address.  On 64-bit systems the odds are better but not zero,
especially for programs that process file offsets, hash values, or
pixel data in the gigabyte range.  Each false positive pins a dead
object (and everything it transitively references) in memory
indefinitely.  Programs that run for a long time with large heaps can
see unbounded memory growth from these phantom roots.

\textbf{No solution within C.}  The C standard deliberately leaves
pointer representation implementation-defined and permits casts
between pointers and integers.  A conforming C program can hide a
pointer in a way that no scanning algorithm can find.  This is not a
bug in the language---it is a consequence of C's design as a portable
assembly language with minimal abstraction over the machine.

Languages like Modula-3 and Java avoid this entirely.  In M3, the
compiler emits stack frame descriptors that tell the garbage collector
exactly which stack slots and object fields contain traced references.
The runtime knows the precise layout of every heap object because
allocations go through \texttt{NEW}, which records the type.  The
collector does not guess: it follows exactly the traced references and
nothing else.  No false negatives, false positives, or heuristics.
Java achieves the same property through the JVM's typed operand stack
and GC maps.  This is what makes the unified-heap architecture of
MScheme, JScheme, and Kawa possible---the collector's precision is
guaranteed by the language, not hoped for by convention.


\subsection{Registration Direction: Push vs.\ Pull}
\label{sec:registration}

\textbf{Push} (the host declares what's available to the script):
MScheme, Tinylisp, Lua, Guile, Tcl, Python.

\textbf{Pull} (the script declares what it wants from the host):
CFFI, Racket FFI.

\textbf{Neither} (shared runtime): Kawa, JScheme.

MScheme is firmly push-based, but its push is \emph{automated} via
\texttt{sstubgen}.  The programmer adds
\texttt{SchemeStubs("InterfaceName")} to the makefile, and the build
system generates and compiles all binding code.  The effect feels
almost like Kawa's shared runtime: everything visible in M3 interfaces
becomes visible in Scheme.


\subsection{Expressiveness at the Boundary}
\label{sec:expressiveness}

MScheme's design allows patterns that are difficult or impossible in
most other systems:

\textbf{M3 objects as first-class Scheme data.}  A Scheme list can
contain M3 \texttt{TextSeq.T} objects alongside Scheme pairs and
numbers.  \texttt{map}, \texttt{filter}, and \texttt{for-each} work
on these mixed lists.  The M3 objects participate in Scheme's
structural equality and serialization (via \texttt{write} and
pickling).

\textbf{Scheme closures as M3 callbacks.}  The \texttt{sstubgen}
documentation describes how M3 object methods can be overridden with
Scheme lambdas, enabling a pattern where M3 objects have individual
methods replaced at runtime---something M3 itself cannot do (methods
are fixed at object allocation time).

\textbf{Runtime discovery.}  Scheme code can enumerate all registered
M3 types, query their supertypes, and list available operations.  This
enables Scheme-level metaprogramming over the M3 type system---writing
generic operations that work on any M3 object based on runtime type
introspection.


%% ==================================================================
\section{Positioning in the Design Space}
\label{sec:positioning}
%% ==================================================================

Plotting the systems on two axes---\emph{per-call cost} (boundary
overhead) vs.\ \emph{setup cost} (how much binding code must be
written)---reveals three clusters:

\begin{samepage}
\begin{verbatim}
                    low setup cost
                         |
            Kawa  *      |
                         |
      MScheme  *         |
    Tinylisp  *          |       * JScheme
  -------------------------------------
  low per-call cost      |       high per-call cost
                         |
                         |      * Guile
              * CFFI     |      * Lua
              * Racket   |      * Tcl
                         |      * Python/C
                         |
                    high setup cost
\end{verbatim}
\end{samepage}

Kawa achieves the ideal: zero setup, zero per-call cost.  But it
requires both languages to run on the JVM.

JScheme shares Kawa's low setup cost (both live on the JVM, no
bindings to write) but pays a significant per-call cost: every Java
method call goes through reflection, while Kawa compiles to direct
\texttt{invokevirtual} bytecode.

MScheme and Tinylisp occupy a similar region: both achieve very low
per-call cost through the shared \texttt{REFANY} heap, and both have
stub generators that reduce setup cost.  MScheme's \texttt{sstubgen}
reduces setup cost further by reading standard \texttt{.i3} files
directly, while Tinylisp's \texttt{compile\_tli} requires hand-written
\texttt{.tli} files.  Tinylisp compensates with lower per-call
overhead at the Lisp level (compiled vs.\ interpreted).

MScheme comes closest to Kawa's position among systems where the host
language has its own native runtime.  The shared GC eliminates
per-call cost for value passing, and \texttt{sstubgen} eliminates
setup cost for binding generation.  The remaining cost is method
dispatch (atom comparison), which is comparable to Lua's metatable
lookup and far cheaper than Python's MRO traversal.

The C-hosted systems (Lua, Guile, Tcl, Python) all pay for the
impedance mismatch between a non-GC host and a GC'd guest.  CFFI and
Racket minimize setup cost on the host side (the C library is
unmodified) but pay for it on the scripting side (manual declarations)
and at every call boundary (marshaling).


%% ==================================================================
\section{Native Code, Not a Virtual Machine}
\label{sec:native-code}
%% ==================================================================

The preceding analysis groups MScheme with JScheme and Kawa as
``shared heap'' systems.  But there is a further distinction that
separates MScheme from its JVM counterparts, and connects it to a much
smaller lineage.


\subsection{Why This Distinction Matters}
\label{sec:why-native}

JScheme and Kawa achieve glueless interoperability on the JVM.  But
the JVM is a virtual machine: Java source compiles to bytecode, which
is then interpreted or JIT-compiled by the VM at runtime.  The GC,
runtime type information, and universal reference type
(\texttt{java.lang.Object}) are services provided by the VM, not by
the compiled program itself.

Modula-3 is different.  The cm3 compiler emits native machine code---x86,
x86-64, ARM64, SPARC, PA-RISC, PowerPC---directly, with no
intermediate bytecode and no virtual machine.  The generated code runs
as ordinary machine instructions.  Yet Modula-3 still provides precise
garbage collection, runtime type discrimination (\texttt{TYPECASE},
\texttt{ISTYPE}, \texttt{TYPECODE}), and a universal traced reference
type (\texttt{REFANY}).  It achieves this by having the compiler emit
metadata alongside the native code: stack frame descriptors that tell
the runtime exactly which stack slots contain traced references, and
type descriptors that record the layout of every heap-allocated type.

This means MScheme's Scheme interpreter runs as native machine code
calling into native machine code, with no VM layer in between.  When
Scheme code calls an M3 method, it is a direct procedure call---the
same calling convention as any M3-to-M3 call.  When the garbage
collector runs, it traces M3 objects and Scheme pairs with the same
algorithm, using the same compiler-generated stack maps.  There is no
bytecode interpreter, JIT warmup, or tiered compilation.


\subsection{How Rare Is This Architecture?}
\label{sec:how-rare}

The combination of properties that makes MScheme possible is unusual:
(1)~ahead-of-time compilation to native code, (2)~precise tracing GC,
(3)~runtime type information with a universal reference type, and
(4)~an embedded interpreter for a \emph{different} language that
shares the host's values and GC directly, with no marshaling.

Very few languages in history have provided (1)--(3) simultaneously.
The lineage leading to Modula-3---Mesa $\to$ Cedar $\to$ Modula-2+
$\to$ Modula-3---is one of the few that does (see
Section~\ref{sec:mesa-cedar}).

A survey of other natively-compiled GC'd languages shows why MScheme's
architecture is hard to replicate:

\textbf{Common Lisp (SBCL, CMUCL, etc.).}  The strongest analogue.
SBCL compiles to native x86-64 and ARM64 code and has a precise
generational GC.  Its interpreter (when enabled) shares the exact same
heap, the same tagged-pointer value representation, and the same
collector as compiled code.  A Scheme interpreter written in Common
Lisp would be architecturally identical to MScheme.  However, CL's
interpreter is self-hosting---it interprets Common Lisp, not a
different language.  The system is designed for mixed-mode execution of
a single language, not for embedding a scripting language in a systems
language.

\textbf{Smalltalk (StrongTalk, Cog/Spur).}  JIT-compiles methods to
native code; the interpreter and compiled code share a single object
space with no marshaling boundary.  Again self-hosting: the
interpreted and compiled languages are the same Smalltalk.
Compilation is JIT (at runtime), not AOT (ahead of time).

\textbf{Eiffel (EiffelStudio's Melting Ice).}  Compiles Eiffel to C
and thence to native code (``frozen'' code), but also maintains
bytecode-interpreted ``melted'' code for rapid development.  Frozen
and melted code share the same heap and objects.  Self-hosting: the
interpreted language is Eiffel.

\textbf{Julia.}  JIT-compiles to native code via LLVM; \texttt{eval}
operates in the same runtime with the same values.  Self-hosting; JIT,
not AOT.

\textbf{Go (Yaegi).}  Compiles to native code and has precise GC, but
Yaegi (a Go interpreter written in Go) wraps values in
\texttt{reflect.Value} at every boundary---a form of marshaling.  Go
also lacks a universal reference type; \texttt{interface\{\}} is a
two-word pair, not a traced pointer.

\textbf{Oberon.}  Compiles to native code and has GC, but lacks a
universal reference type and has no embedded interpreter---it uses
dynamic loading of compiled modules instead.

\textbf{D, Dylan, OCaml.}  All compile to native code and have GC,
but none has an embedded interpreter for a different language that
shares the host heap without marshaling.  OCaml's bytecode toplevel
shares the heap with loaded \texttt{.cmo} modules, but the
native-code compiler path has no analogous mechanism for a different
language.

The pattern that emerges is that self-hosting interpreters in
natively-compiled languages are not uncommon (CL, Smalltalk, Eiffel,
Julia all have them), but an interpreter for a \emph{different}
language that shares the host's heap with zero marshaling overhead is
rare.  The obstacle is not theoretical---any language with GC, RTTI,
and a universal reference type could support it---but practical: very
few such languages exist, and in those that do, the idea has seldom
been pursued.  The one important exception---Tinylisp
(Section~\ref{sec:tinylisp})---exploits the same mechanism in
Modula-2+, the direct ancestor of Modula-3.


%% ==================================================================
\section{Design Philosophy: Value Types and the Lisp Machine}
\label{sec:philosophy}
%% ==================================================================

The SSTUBGEN.TXT design document explicitly positions the system's
goal as making Scheme a viable ``main program'' language for large M3
applications: the libraries are in M3, the orchestration is in Scheme.
This is the same pattern later popularized by Lua in game engines and
Python in scientific computing---but MScheme achieves it with
significantly less impedance mismatch, thanks to the shared GC and
type system.

Why do so few languages provide the features that make this
architecture possible?  The answer leads to a deeper question about
language design and the relationship between MScheme's architecture
and the Lisp Machine.

The Lisp Machine is the ultimate precedent for glueless
interoperability: Lisp \emph{was} the systems language and no
interoperability boundary existed at all.  But the Lisp Machine also
illustrates the cost of the approach the Algol family avoids.  As
Nelson observes in \emph{Systems Programming with
Modula-3}~\cite{nelson1991}:

\begin{quote}
At the other extreme are languages like Lisp, ML, Smalltalk, and CLU,
whose programming models originate from mathematics.\ [\ldots]\ These
languages have beautiful programming models, but they tend to be
difficult to implement efficiently, because the uniform treatment of
values in the programming model invites a runtime system in which
values are uniformly represented by pointers.  If the implementer
doesn't take steps to avoid it, as simple a statement as
\texttt{n := n + 1} could require an allocation, a method lookup, or
both.
\end{quote}

The Lisp Machine solved this in hardware.  The Symbolics 3600
(Moon~\cite{moon1985}) and its successor the Ivory (I-Machine) use
tagged words: every machine word carries a data type tag (4 bits on
the 3600, 6 bits on the Ivory) alongside 32 bits of data or address.
Twenty-two of the Ivory's type codes designate \emph{immediate}
values---fixnums, single-precision floats, characters, small
ratios---where the data is encoded directly in the word.  Fixnum
arithmetic is a register-to-register operation; \texttt{n + 1}
requires no allocation and no GC involvement.  The remaining 33 type
codes designate \emph{pointer} values---lists, arrays, symbols,
closures, bignums---where the 32-bit field is a heap address.
CDR-coding compresses list cells to roughly half the space of naive
linked lists.  Hardware-assisted ephemeral garbage collection reduces
the cost of short-lived heap allocations.  Common Lisp's
\texttt{dynamic-extent} declaration allows the compiler to
stack-allocate objects known not to escape their enclosing scope.  With
all of these mechanisms, Lisp Machines do handle fixnum-intensive code
efficiently---but they achieve it through \emph{dedicated hardware}
that conventional machines lack.

Tagged-pointer hardware has a longer history---the Burroughs B5000
(1961) and B6500/B6700 (1969)~\cite{organick1973} tagged every word
for machine-level integrity---but the Lisp Machines were the first to
tag by \emph{user-level data type} (fixnum, cons, symbol, closure),
enabling the hardware to dispatch on language-level types directly.

The Algol family, to which Modula-3 belongs, takes the opposite
approach: instead of solving the problem in hardware, the \emph{language
itself} distinguishes value types from reference types.  In Modula-3,
\texttt{INTEGER}, \texttt{REAL}, \texttt{BOOLEAN}, \texttt{CHAR},
records, and fixed arrays are value types---stack-allocated, unboxed,
and invisible to the garbage collector.  Only reference types
(\texttt{REF}, \texttt{OBJECT}, \texttt{TEXT}) participate in the
traced heap.  This means that most of a typical M3 systems
program---local variables, loop counters, parameter passing, record
manipulation---never touches the GC at all.  The garbage collector
runs only for heap-allocated reference types, and the compiler needs
no special optimization or declarations to achieve this: the type
system guarantees it statically.

This is exactly the property that makes MScheme's architecture
practical for systems programming.  The Scheme interpreter and its
data structures (pairs, closures, environments) live in the traced
heap and are managed by the GC, but the M3 code that the Scheme stubs
call into can be dominated by stack-allocated value types.  The
boundary between ``GC'd scripting world'' and ``stack-allocated systems
world'' is not a marshaling barrier---it is the ordinary Modula-3 type
system, which \texttt{sstubgen} navigates automatically by generating
\texttt{ToScheme}/\texttt{ToModula} converters that box and unbox
value types at the boundary while passing reference types through
unchanged.  (JScheme and Kawa on the JVM face a version of the same
issue---the JVM boxes primitives like \texttt{int} to
\texttt{Integer}---but Java's autoboxing covers only a fixed set of
primitive types, while \texttt{sstubgen} handles the full M3 type
system including user-defined records, enumerations, subranges, and
multidimensional arrays.)

The Lisp Machine achieved glueless interoperability by making Lisp the
only language.  JScheme approximated it for Java on the JVM.  MScheme
approximates it for Modula-3 on conventional hardware---one type alias
(\texttt{T = REFANY}), a language that distinguishes values from
references, and a code generator that handles the rest.


%% ==================================================================
\section{Historical Lineage}
\label{sec:lineage}
%% ==================================================================

\subsection{The Mesa/Cedar/Modula-2+/Modula-3 Language Family}
\label{sec:mesa-cedar}

The language features that make MScheme and Tinylisp possible---\texttt{REFANY},
precise tracing GC, \texttt{TYPECASE}---trace a specific lineage
through four decades of systems programming language design.

\textbf{Mesa} (Xerox PARC, 1970s) was a systems language for the Alto
and later the Dorado.  It compiled to native code and had strong
typing, but no garbage collection and no universal reference type.

\textbf{Cedar} (Xerox PARC, early 1980s) extended Mesa with garbage
collection, \texttt{REF ANY}, and runtime type
information~\cite{rovner1985}.  Cedar had the infrastructure to
support an MScheme-like system---a universal traced reference type and
a precise GC---but no one built a full scripting interpreter on top of
it.  The debugger could evaluate expressions in a stopped frame, but
that is not the same as a general-purpose embedded language.

\textbf{Modula-2+} (DEC Systems Research Center, mid-1980s) was DEC
SRC's dialect of Modula-2, adding garbage collection, exception
handling, opaque types, \texttt{REFANY}, and \texttt{TYPECASE}---the
features Cedar had pioneered, adapted to a Wirth-family language.  It
was the implementation language for SRC's systems infrastructure and
the language in which Tinylisp was written.

\textbf{Modula-3} (DEC SRC / Olivetti Research Center,
1989)~\cite{cardelli1989report} refined Modula-2+ into a clean design
with a formal specification.  The Modula-3 Report codified the type
system, the \texttt{REFANY} mechanism, and the \texttt{TYPECASE}
dispatch that make MScheme's architecture possible.


\subsection{From JScheme to MScheme}
\label{sec:jscheme-to-mscheme}

Peter Norvig published JScheme in 1998~\cite{norvig1998} as a compact
demonstration that a useful Scheme interpreter could be written in
\textasciitilde1700 lines of Java.  The key architectural insight---\texttt{java.lang.Object}
as the universal Scheme value---was implicit in the choice of Java as
the implementation language.  Ken Anderson (BBN) and Tim Hickey
(Brandeis) later expanded JScheme~\cite{anderson2003} with
dot-notation syntax, overload resolution, and dynamic proxy support
for Java interfaces.

MScheme was written in 2007--2008 as a line-for-line translation of
JScheme into Modula-3.  The translation preserved JScheme's
fundamental principle (\texttt{Object} $\to$ \texttt{REFANY}) but
added two things JScheme lacks: an automatic stub generator
(\texttt{sstubgen}) and a runtime type introspection API.  MScheme has
been used in financial trading, scientific computing, and hardware
design.

The MScheme design predates many of the modern FFI systems (Racket's
FFI was redesigned around 2010; Python's stable ABI came in 3.2/2011;
Guile's foreign object types replaced SMOBs in 2.0/2011).

MScheme was not aware of Tinylisp; the parallel development reflects
the fact that the architecture is essentially forced once the host
language provides \texttt{REFANY} and GC.


\subsection{The sstubgen and m3tk Provenance}
\label{sec:sstubgen-provenance}

\texttt{sstubgen} derives directly from the Network Objects stub
generator (\texttt{m3-comm/stubgen}) designed by Susan Owicki and Ted
Wobber for Birrell et al.'s distributed object
system~\cite{birrell1994}.  It is not merely modeled on the network
objects \texttt{stubgen}---it shares code.  Several source files are
identical between the two packages (\texttt{FRefRefTbl},
\texttt{TypeNames}, \texttt{ValueProc}), and others
(\texttt{AstToType}, \texttt{Type}, \texttt{Value},
\texttt{StubGenTool}, \texttt{StubUtils}) are extended copies.

M3tk is a comprehensive Modula-3 front-end toolkit providing a generic
AST framework, parser, semantic analyzer, and support for building
compiler extensions.  It was written by Mick Jordan, originally at the
Olivetti Research Center (ORC) in Menlo Park, California, and later at
DEC SRC, and described in his 1990 paper ``An Extensible Programming
Environment for Modula-3''~\cite{jordan1990}.  M3tk began life as part
of Olivetti's full Modula-3 compiler; the cm3 compiler (Bill Kalsow's
SRC compiler lineage) does not use m3tk itself, but the Network
Objects \texttt{stubgen} does, and \texttt{sstubgen} inherits that
dependency.  M3tk is organized into sub-packages (\texttt{m3tk-ast},
\texttt{m3tk-syn}, \texttt{m3tk-sem}, \texttt{m3tk-fe}, etc.)\ and is
used by both \texttt{stubgen} and \texttt{sstubgen}, as well as by
the language server and other M3 development tools.


%% ==================================================================
\section*{Acknowledgments}
%% ==================================================================

Paul McJones provided detailed comments on an earlier draft, including
the reference to Ellis's Tinylisp and corrections to the discussion
of m3tk, stub naming conventions, and the role of the \texttt{interp}
and \texttt{excH} parameters in generated code.


%% ==================================================================
\begin{thebibliography}{99}
%% ==================================================================

%% --- Modula-3 ---

\bibitem{cardelli1989report}
L.~Cardelli, J.~Donahue, L.~Glassman, M.~Jordan, B.~Kalsow, and
G.~Nelson.
\newblock Modula-3 Report (revised).
\newblock DEC Systems Research Center, Report~52, November 1989.

\bibitem{cardelli1989types}
L.~Cardelli, J.~Donahue, M.~Jordan, B.~Kalsow, and G.~Nelson.
\newblock The {Modula-3} type system.
\newblock In \emph{Conference Record of the Sixteenth Annual ACM
  Symposium on Principles of Programming Languages}, pp.~202--212,
  1989.

\bibitem{cardelli1989typeful}
L.~Cardelli.
\newblock Typeful programming.
\newblock DEC Systems Research Center, Report~45, May 1989.
\newblock Revised version in E.\,J.\ Neuhold and M.\ Paul (eds.),
  \emph{Formal Description of Programming Concepts},
  Springer-Verlag, 1991.

\bibitem{nelson1991}
G.~Nelson (ed.).
\newblock \emph{Systems Programming with Modula-3.}
\newblock Prentice Hall, 1991.

\bibitem{harbison1992}
S.\,P.\ Harbison.
\newblock \emph{Modula-3.}
\newblock Prentice Hall, 1992.

\bibitem{cm3}
Modula3/CM3: Critical Mass Modula-3.
\newblock \url{https://github.com/modula3/cm3}

%% --- Network Objects and stub generation ---

\bibitem{birrell1994}
A.\,D.\ Birrell, G.~Nelson, S.~Owicki, and E.~Wobber.
\newblock Network objects.
\newblock DEC Systems Research Center, Report~115, February 1994.
\newblock Also in \emph{Proc.\ 14th ACM SOSP}, December 1993; and
  \emph{Software---Practice and Experience},
  25(S4):87--130, December 1995.

\bibitem{birrell1993gc}
A.\,D.\ Birrell, D.~Evers, G.~Nelson, S.~Owicki, and E.~Wobber.
\newblock Distributed garbage collection for network objects.
\newblock DEC Systems Research Center, Report~116, December 1993.

%% --- M3 Toolkit ---

\bibitem{jordan1990}
M.~Jordan.
\newblock An extensible programming environment for {Modula-3}.
\newblock In \emph{Proc.\ Fourth ACM SIGSOFT Symposium on Software
  Development Environments} (SDE~4), Irvine, California, December
  1990.
\newblock \emph{ACM SIGSOFT Software Engineering Notes},
  15(6):66--76.

\bibitem{jordan1988}
M.~Jordan and P.~Robinson.
\newblock A programming environment for {Modula-2}.
\newblock \emph{Software Engineering Journal}, 3(4):119--126, July
  1988.

%% --- MScheme, sstubgen, cspc ---

\bibitem{mscheme}
MScheme source code: \texttt{cm3/m3-scheme/mscheme/}.
\newblock In the CM3 repository.

\bibitem{sstubgen-txt}
SSTUBGEN.TXT: Design document for the Scheme stub generator.
\newblock \texttt{cm3/m3-scheme/sstubgen/src/SSTUBGEN.TXT}.

\bibitem{schemem3-readme}
SchemeM3 README: Embedding guide.
\newblock \texttt{cm3/m3-scheme/modula3scheme/src/README}.

\bibitem{cspc}
\texttt{cspc}: CSP compiler for asynchronous VLSI design.
\newblock \texttt{intel-async/async-toolkit/m3utils/m3utils/csp/src/}.
\newblock In the async-toolkit repository.

%% --- JScheme ---

\bibitem{norvig1998}
P.~Norvig.
\newblock {JScheme}: {Scheme} implemented in {Java}.
\newblock \url{https://norvig.com/jscheme.html}, 1998.

\bibitem{norvig-design}
P.~Norvig.
\newblock {JScheme}: Design decisions.
\newblock \url{https://norvig.com/jscheme-design.html}

\bibitem{anderson2003}
K.\,R.\ Anderson, T.\,J.\ Hickey, and P.~Norvig.
\newblock {JScheme}.
\newblock \url{https://www.cs.brandeis.edu/~tim/jscheme/main.html}.
\newblock Brandeis University / BBN Technologies, 2003.

%% --- Tinylisp and Modula-2+ ---

\bibitem{ellis1989}
J.~Ellis.
\newblock Tinylisp reference manual.
\newblock Systems Research Center, Digital Equipment Corporation,
  January 26, 1989.
\newblock \url{https://softwarepreservation.computerhistory.org/LISP/dec/Ellis-Tinylisp-1989.pdf}

\bibitem{m2plus}
Modula-2+ manuals and related materials.
\newblock \url{https://softwarepreservation.computerhistory.org/modula2+/}

%% --- Cedar and Modula-3 ancestry ---

\bibitem{rovner1985}
P.~Rovner.
\newblock On adding garbage collection and runtime types to a
  strongly-typed, statically-checked, concurrent language.
\newblock Xerox PARC, CSL-84-7, July 1985.

\bibitem{lampson1983}
B.\,W.\ Lampson.
\newblock A description of the {Cedar} language.
\newblock Xerox PARC, CSL-83-15, December 1983.

%% --- Lua ---

\bibitem{lua-manual}
R.~Ierusalimschy, L.\,H.\ de~Figueiredo, and W.~Celes.
\newblock \emph{Lua 5.4 Reference Manual.}
\newblock \url{https://www.lua.org/manual/5.4/}

\bibitem{lua-pil}
R.~Ierusalimschy.
\newblock \emph{Programming in Lua}, 4th edition.
\newblock Lua.org, 2016.

\bibitem{lua-evolution}
R.~Ierusalimschy, L.\,H.\ de~Figueiredo, and W.~Celes.
\newblock The evolution of {Lua}.
\newblock In \emph{Proc.\ Third ACM SIGPLAN Conference on History of
  Programming Languages} (HOPL~III), 2007.

%% --- GNU Guile ---

\bibitem{guile-manual}
\emph{GNU Guile Reference Manual}, version~3.0.
\newblock \url{https://www.gnu.org/software/guile/manual/}

%% --- Tcl/Tk ---

\bibitem{ousterhout1994}
J.\,K.\ Ousterhout.
\newblock \emph{Tcl and the Tk Toolkit.}
\newblock Addison-Wesley, 1994.

\bibitem{tcl-doc}
\emph{Tcl/Tk Documentation.}
\newblock \url{https://www.tcl-lang.org/doc/}

%% --- CPython ---

\bibitem{python-capi}
\emph{Python/C API Reference Manual.}
\newblock \url{https://docs.python.org/3/c-api/}

\bibitem{python-extending}
\emph{Extending and Embedding the Python Interpreter.}
\newblock \url{https://docs.python.org/3/extending/}

%% --- Kawa Scheme ---

\bibitem{bothner1998}
P.~Bothner.
\newblock {Kawa}---compiling dynamic languages to the {Java VM}.
\newblock In \emph{Proc.\ USENIX Annual Technical Conference},
  FREENIX Track, 1998.

\bibitem{kawa-manual}
\emph{The Kawa Scheme Language.}
\newblock \url{https://www.gnu.org/software/kawa/}

%% --- Common Lisp CFFI ---

\bibitem{cffi}
\emph{CFFI: The Common Foreign Function Interface.}
\newblock \url{https://cffi.common-lisp.dev/}

\bibitem{cffi-manual}
\emph{CFFI User Manual.}
\newblock \url{https://cffi.common-lisp.dev/manual/cffi-manual.html}

%% --- Racket FFI ---

\bibitem{barzilay2004}
E.~Barzilay and D.~Orlovsky.
\newblock Foreign interface for {PLT Scheme}.
\newblock In \emph{Proc.\ Fifth ACM SIGPLAN Workshop on Scheme and
  Functional Programming}, 2004.

\bibitem{racket-ffi}
\emph{The Racket Foreign Interface.}
\newblock \url{https://docs.racket-lang.org/foreign/}

%% --- Conservative garbage collection ---

\bibitem{boehm1988}
H.-J.\ Boehm and M.~Weiser.
\newblock Garbage collection in an uncooperative environment.
\newblock \emph{Software---Practice and Experience},
  18(9):807--820, September 1988.

\bibitem{boehm1991}
H.-J.\ Boehm, A.\,J.\ Demers, and S.~Shenker.
\newblock Mostly parallel garbage collection.
\newblock In \emph{Proc.\ ACM SIGPLAN '91 Conference on Programming
  Language Design and Implementation}, pp.~157--164, 1991.

%% --- SBCL ---

\bibitem{sbcl-manual}
\emph{SBCL User Manual.}
\newblock \url{https://www.sbcl.org/manual/}

\bibitem{sbcl-source}
Steel Bank Common Lisp source.
\newblock \url{https://github.com/sbcl/sbcl}

%% --- Eiffel ---

\bibitem{meyer1997}
B.~Meyer.
\newblock \emph{Object-Oriented Software Construction}, 2nd edition.
\newblock Prentice Hall, 1997.

\bibitem{melting-ice}
Melting Ice Technology.
\newblock \url{https://www.eiffel.org/doc/eiffelstudio/Melting_Ice_Technology}

%% --- Smalltalk and StrongTalk ---

\bibitem{bracha1993}
G.~Bracha and D.~Griswold.
\newblock Strongtalk: Typechecking {Smalltalk} in a production
  environment.
\newblock In \emph{Proc.\ ACM Conference on Object-Oriented
  Programming, Systems, Languages and Applications} (OOPSLA),
  pp.~215--230, 1993.

\bibitem{ungar1987}
D.~Ungar and R.\,B.\ Smith.
\newblock Self: The power of simplicity.
\newblock In \emph{Proc.\ OOPSLA~'87}, pp.~227--241, 1987.

%% --- Oberon ---

\bibitem{wirth1992}
N.~Wirth and J.~Gutknecht.
\newblock \emph{Project Oberon: The Design of an Operating System and
  Compiler.}
\newblock ACM Press / Addison-Wesley, 1992.
\newblock Revised edition (\emph{The Design of an Operating System, a
  Compiler, and a Computer}), 2013.

%% --- Lisp Machines and tagged architectures ---

\bibitem{moon1985}
D.\,A.\ Moon.
\newblock Architecture of the {Symbolics}~3600.
\newblock In \emph{Proc.\ 12th Annual International Symposium on
  Computer Architecture} (ISCA), pp.~76--83, 1985.

\bibitem{organick1973}
E.\,I.\ Organick.
\newblock \emph{Computer System Organization: The B5700/B6700 Series.}
\newblock Academic Press, 1973.

\bibitem{levy1984}
H.\,M.\ Levy.
\newblock \emph{Capability-Based Computer Systems.}
\newblock Digital Press, 1984.
\newblock Chapter~2 covers the Burroughs B5000 and B6500 tagged
  architectures.

%% --- Other systems cited ---

\bibitem{yaegi}
Traefik.
\newblock Yaegi: Another elegant {Go} interpreter.
\newblock \url{https://github.com/traefik/yaegi}

\bibitem{hint}
Hint: Runtime {Haskell} interpreter.
\newblock \url{https://hackage.haskell.org/package/hint}

\bibitem{racket-guide}
M.~Flatt, R.\,B.\ Findler, and PLT.
\newblock \emph{The Racket Guide.}
\newblock \url{https://docs.racket-lang.org/guide/}

\end{thebibliography}

\end{document}
